% !TEX spellcheck = en_US



% ~~~~~~~~~~~~~~~~~~~~~~~~~~~~~~~~~~~~~~~ V
% (NC) 2026 Haluk Bingol
% github.com/halukbingol/LaTeX-Templates
%
% Licensees may copy, distribute, display, and perform the work and make derivative works 
% and remixes based on it only for non-commercial purposes. 
% ~~~~~~~~~~~~~~~~~~~~~~~~~~~~~~~~~~~~~~~ A


% v20261029T1800

\documentclass[aspectratio=169,10pt,compress]{beamer}




% ~~~~~~~~~~~~~~~~~~~~~~~~~~~~~~~~~~~~~~~~~~~~~~~~~~~~~~~~~~~~~~~~~ V hblib
	\usepackage{../../hblib/hblibPresentation} % lib presentation: theme, i/N, \hSection, algorithm2e, ...
	\usepackage{../../hblib/hblibCodeoutput} % lib code + output: \lstcode, \lstterminal, ...
	\usepackage{../../hblib/hblibDatastructures} % lib data structures: \hString, \hStack, \hQueue, \hArray
% ~~~~~~~~~~~~~~~~~~~~~~~~~~~~~~~~~~~~~~~~~~~~~~~~~~~~~~~~~~~~~~~~~ A hblib




% +++++++++++++++++++++++++++++++++++++++ V specific
\DontPrintSemicolon % in algorithms
% +++++++++++++++++++++++++++++++++++++++ A specific



\title[Arrays 1D: Examples]{
	One-Dimensional Arrays\\
	Examples
}
\subtitle{
	\cisOOO\ -- Level 5: every concept as a small program
}
\author[Erkovan]{Sait Erkovan}

\institute[FCIS]{
	Faculty of Computer and Information Sciences\\
	Yeditepe University
}
\date{{\tiny \hbTimeStamp}}

\begin{document}




% ~~~~~~~~~~~~~~~~~~~~~~~~~~~~~~~~~~~~~~ V frm
\begin{frame}[t,fragile]
	\titlepage
\end{frame}
% ~~~~~~~~~~~~~~~~~~~~~~~~~~~~~~~~~~~~~~ A frm




% ~~~~~~~~~~~~~~~~~~~~~~~~~~~~~~~~~~~~~~ V frm
\begin{frame}[t,fragile] \frametitle{Outline}
	\begin{columns}[T]
		\begin{column}{0.48\textwidth}
			\tableofcontents[sections={1-4}, hideallsubsections]
		\end{column}
		\begin{column}{0.48\textwidth}
			\tableofcontents[sections={5-8}, hideallsubsections]
		\end{column}
	\end{columns}

	\vspace{1em}
	\small Each slide shows a program first. \hQuestion{Predict its output}, then go to the next
	click. The sections follow \hCode{array1D-0-concepts}; the pictures are there.
	Every output here was produced by \hCode{run\_examples.sh}.
\end{frame}
% ~~~~~~~~~~~~~~~~~~~~~~~~~~~~~~~~~~~~~~ A frm




% ~~~~~~~~~~~~~~~~~~~~~~~~~~~~~~~~~~~~~~ V
\hSection[Why]{Why arrays?}
% ~~~~~~~~~~~~~~~~~~~~~~~~~~~~~~~~~~~~~~ A




% ~~~~~~~~~~~~~~~~~~~~~~~~~~~~~~~~~~~~~~ V frm
\begin{frame}[t,fragile] \frametitle{Days warmer than the average}
	\lstcoderightstep[0.66][basicstyle=\ttfamily\scriptsize]{codeoutput/AboveAverage.java}{Java}{codeoutput/AboveAverage.txt}

	\footnotesize Concept: some questions need every value \hEmph{twice}: once for the average,
	once to compare.
\end{frame}
% ~~~~~~~~~~~~~~~~~~~~~~~~~~~~~~~~~~~~~~ A frm




% ~~~~~~~~~~~~~~~~~~~~~~~~~~~~~~~~~~~~~~ V
\hSection[Anatomy]{Anatomy of an array}
% ~~~~~~~~~~~~~~~~~~~~~~~~~~~~~~~~~~~~~~ A




% ~~~~~~~~~~~~~~~~~~~~~~~~~~~~~~~~~~~~~~ V frm
\begin{frame}[t,fragile] \frametitle{Length, first index, last index}
	\lstcoderightstep[0.66][basicstyle=\ttfamily\scriptsize]{codeoutput/Anatomy.java}{Java}{codeoutput/Anatomy.txt}

	\small Concept: indices run from 0 to \hCode{length - 1}; the length is a count.
\end{frame}
% ~~~~~~~~~~~~~~~~~~~~~~~~~~~~~~~~~~~~~~ A frm




% ~~~~~~~~~~~~~~~~~~~~~~~~~~~~~~~~~~~~~~ V
\hSection[Rules]{The rules of a Java array}
% ~~~~~~~~~~~~~~~~~~~~~~~~~~~~~~~~~~~~~~ A




% ~~~~~~~~~~~~~~~~~~~~~~~~~~~~~~~~~~~~~~ V frm
\begin{frame}[t,fragile] \frametitle{Default values}
	\lstcoderightstep[0.66][basicstyle=\ttfamily\scriptsize]{codeoutput/Defaults.java}{Java}{codeoutput/Defaults.txt}

	\small Concept: a new array is never empty; every element starts with the default of its type.
\end{frame}
% ~~~~~~~~~~~~~~~~~~~~~~~~~~~~~~~~~~~~~~ A frm




% ~~~~~~~~~~~~~~~~~~~~~~~~~~~~~~~~~~~~~~ V
\hSection[Memory]{Arrays live in memory}
% ~~~~~~~~~~~~~~~~~~~~~~~~~~~~~~~~~~~~~~ A




% ~~~~~~~~~~~~~~~~~~~~~~~~~~~~~~~~~~~~~~ V frm
\begin{frame}[t,fragile] \frametitle{Alias or copy?}
	\lstcoderightstep[0.66][basicstyle=\ttfamily\scriptsize]{codeoutput/Aliasing.java}{Java}{codeoutput/Aliasing.txt}

	\small Concept: \hCode{b = a} copies the \emph{arrow}; a copy is a second array.
	\hCode{==} asks ``same array?''.
\end{frame}
% ~~~~~~~~~~~~~~~~~~~~~~~~~~~~~~~~~~~~~~ A frm




% ~~~~~~~~~~~~~~~~~~~~~~~~~~~~~~~~~~~~~~ V
\hSection[Patterns]{Thinking in traversals}
% ~~~~~~~~~~~~~~~~~~~~~~~~~~~~~~~~~~~~~~ A




% ~~~~~~~~~~~~~~~~~~~~~~~~~~~~~~~~~~~~~~ V frm
\begin{frame}[t,fragile] \frametitle{Fill: the value depends on the index}
	\lstcoderightstep[0.66][basicstyle=\ttfamily\scriptsize]{codeoutput/FillSquares.java}{Java}{codeoutput/FillSquares.txt}

	\small Pattern \emph{fill}: $a[i] \gets f(i)$; here $f(i) = i^{2}$.
\end{frame}
% ~~~~~~~~~~~~~~~~~~~~~~~~~~~~~~~~~~~~~~ A frm




% ~~~~~~~~~~~~~~~~~~~~~~~~~~~~~~~~~~~~~~ V frm
\begin{frame}[t,fragile] \frametitle{Accumulate: sum and average}
	\lstcoderightstep[0.66][basicstyle=\ttfamily\scriptsize]{codeoutput/SumAverage.java}{Java}{codeoutput/SumAverage.txt}

	\small Pattern \emph{accumulate}. \hQuestion{Why \hCode{(double)}?}
	\uncover<2->{Without it, \hCode{389 / 5} is 77. And $0.0 / 0$ is \hCode{NaN}, ``not a number''.}
\end{frame}
% ~~~~~~~~~~~~~~~~~~~~~~~~~~~~~~~~~~~~~~ A frm




% ~~~~~~~~~~~~~~~~~~~~~~~~~~~~~~~~~~~~~~ V frm
\begin{frame}[t,fragile] \frametitle{Best so far: wettest and driest}
	\lstcoderightstep[0.66][basicstyle=\ttfamily\scriptsize]{codeoutput/MinMax.java}{Java}{codeoutput/MinMax.txt}

	\small Pattern \emph{best so far}, twice in one loop. \hQuestion{Wed and Fri tie. Why Wed?}
\end{frame}
% ~~~~~~~~~~~~~~~~~~~~~~~~~~~~~~~~~~~~~~ A frm




% ~~~~~~~~~~~~~~~~~~~~~~~~~~~~~~~~~~~~~~ V frm
\begin{frame}[t,fragile] \frametitle{Search: the first match, or -1}
	\lstcoderightstep[0.66][basicstyle=\ttfamily\scriptsize]{codeoutput/LinearSearch.java}{Java}{codeoutput/LinearSearch.txt}

	\small Pattern \emph{search}, written with \hCode{return} inside the loop: the method ends at the first match.
\end{frame}
% ~~~~~~~~~~~~~~~~~~~~~~~~~~~~~~~~~~~~~~ A frm




% ~~~~~~~~~~~~~~~~~~~~~~~~~~~~~~~~~~~~~~ V frm
\begin{frame}[t,fragile] \frametitle{Count the matches}
	\lstcoderightstep[0.66][basicstyle=\ttfamily\scriptsize]{codeoutput/CountPassing.java}{Java}{codeoutput/CountPassing.txt}

	\small Pattern \emph{count}: a for-each loop is enough, because the index is not needed.
\end{frame}
% ~~~~~~~~~~~~~~~~~~~~~~~~~~~~~~~~~~~~~~ A frm




% ~~~~~~~~~~~~~~~~~~~~~~~~~~~~~~~~~~~~~~ V frm
\begin{frame}[t,fragile] \frametitle{Transform in place}
	\lstcoderightstep[0.66][basicstyle=\ttfamily\scriptsize]{codeoutput/CurveScores.java}{Java}{codeoutput/CurveScores.txt}

	\small Pattern \emph{transform}: writing \hCode{scores[i]} needs the index loop.
\end{frame}
% ~~~~~~~~~~~~~~~~~~~~~~~~~~~~~~~~~~~~~~ A frm




% ~~~~~~~~~~~~~~~~~~~~~~~~~~~~~~~~~~~~~~ V frm
\begin{frame}[t,fragile] \frametitle{Filter into a new array}
	\lstcoderightstep[0.66][basicstyle=\ttfamily\scriptsize]{codeoutput/FilterPassing.java}{Java}{codeoutput/FilterPassing.txt}
\end{frame}
% ~~~~~~~~~~~~~~~~~~~~~~~~~~~~~~~~~~~~~~ A frm




% ~~~~~~~~~~~~~~~~~~~~~~~~~~~~~~~~~~~~~~ V frm
\begin{frame}[t,fragile] \frametitle{Swap: wrong and right}
	\lstcoderightstep[0.66][basicstyle=\ttfamily\scriptsize]{codeoutput/Swap.java}{Java}{codeoutput/Swap.txt}

	\small The first swap loses 5 forever. \hCode{temp} keeps it until it is needed.
\end{frame}
% ~~~~~~~~~~~~~~~~~~~~~~~~~~~~~~~~~~~~~~ A frm




% ~~~~~~~~~~~~~~~~~~~~~~~~~~~~~~~~~~~~~~ V frm
\begin{frame}[t,fragile] \frametitle{Fence posts: separators}
	\lstcoderightstep[0.66][basicstyle=\ttfamily\scriptsize]{codeoutput/Separators.java}{Java}{codeoutput/Separators.txt}

	\small A separator \emph{before} every element except the first.
\end{frame}
% ~~~~~~~~~~~~~~~~~~~~~~~~~~~~~~~~~~~~~~ A frm




% ~~~~~~~~~~~~~~~~~~~~~~~~~~~~~~~~~~~~~~ V frm
\begin{frame}[t,fragile] \frametitle{Two patterns in one loop}
	\lstcoderightstep[0.66][basicstyle=\ttfamily\scriptsize]{codeoutput/DropLowest.java}{Java}{codeoutput/DropLowest.txt}

	\small Accumulate and best so far together: the average of the quizzes without the lowest one.
\end{frame}
% ~~~~~~~~~~~~~~~~~~~~~~~~~~~~~~~~~~~~~~ A frm




% ~~~~~~~~~~~~~~~~~~~~~~~~~~~~~~~~~~~~~~ V
\hSection[Partial]{Partially filled arrays}
% ~~~~~~~~~~~~~~~~~~~~~~~~~~~~~~~~~~~~~~ A




% ~~~~~~~~~~~~~~~~~~~~~~~~~~~~~~~~~~~~~~ V frm
\begin{frame}[t,fragile] \frametitle{Size travels with the array}
	\lstcode[basicstyle=\ttfamily\scriptsize, lastline=16]{codeoutput/PartiallyFilled.java}{Java}

	\small \hCode{show} loops to \hCode{size}, not to \hCode{a.length}: the gray cells of the
	concepts deck are never printed.
\end{frame}
% ~~~~~~~~~~~~~~~~~~~~~~~~~~~~~~~~~~~~~~ A frm




% ~~~~~~~~~~~~~~~~~~~~~~~~~~~~~~~~~~~~~~ V frm
\begin{frame}[t,fragile] \frametitle{Insert and remove}
	\lstcoderightstep[0.66][basicstyle=\ttfamily\scriptsize, firstline=18, firstnumber=18]{codeoutput/PartiallyFilled.java}{Java}{codeoutput/PartiallyFilled.txt}
\end{frame}
% ~~~~~~~~~~~~~~~~~~~~~~~~~~~~~~~~~~~~~~ A frm




% ~~~~~~~~~~~~~~~~~~~~~~~~~~~~~~~~~~~~~~ V
\hSection[Methods]{Arrays and methods}
% ~~~~~~~~~~~~~~~~~~~~~~~~~~~~~~~~~~~~~~ A




% ~~~~~~~~~~~~~~~~~~~~~~~~~~~~~~~~~~~~~~ V frm
\begin{frame}[t,fragile] \frametitle{Three methods that take an array}
	\lstcode[basicstyle=\ttfamily\scriptsize, lastline=18]{codeoutput/PassReference.java}{Java}

	\small \hQuestion{Which of them can change the caller's array?}
\end{frame}
% ~~~~~~~~~~~~~~~~~~~~~~~~~~~~~~~~~~~~~~ A frm




% ~~~~~~~~~~~~~~~~~~~~~~~~~~~~~~~~~~~~~~ V frm
\begin{frame}[t,fragile] \frametitle{Calling them}
	\lstcoderightstep[0.66][basicstyle=\ttfamily\scriptsize, firstline=20, firstnumber=20]{codeoutput/PassReference.java}{Java}{codeoutput/PassReference.txt}

	\begin{itemize}\small
		\item<2-> \hCode{addBonus} changes the elements: the caller sees 75 85 95
		\item<2-> \hCode{replace} moves only its own arrow: nothing changes for the caller
		\item<2-> \hCode{doubled} leaves \hCode{s} alone and returns a new array
	\end{itemize}
\end{frame}
% ~~~~~~~~~~~~~~~~~~~~~~~~~~~~~~~~~~~~~~ A frm




% ~~~~~~~~~~~~~~~~~~~~~~~~~~~~~~~~~~~~~~ V
\hSection[Summary]{Summary}
% ~~~~~~~~~~~~~~~~~~~~~~~~~~~~~~~~~~~~~~ A




% ~~~~~~~~~~~~~~~~~~~~~~~~~~~~~~~~~~~~~~ V frm
\begin{frame}[t,fragile] \frametitle{Recipes: concept and program}


% +++++++++++++++++++++++++++++++++++++++ V tbl
%: -Tbl. recipes
\begin{table}
		\footnotesize
		\begin{tabular}{@{}L{0.24\textwidth}L{0.24\textwidth}L{0.44\textwidth}@{}}
			\toprule
			\textbf{Concept} & \textbf{Program} & \textbf{The line to remember} \\
			\midrule
			keep values for later  & \hCode{AboveAverage}    & two passes over the same array \\
			index and length       & \hCode{Anatomy}         & \hCode{a[a.length - 1]} is the last one \\
			default values         & \hCode{Defaults}        & \hCode{0}, \hCode{0.0}, \hCode{false}, code 0, \hCode{null} \\
			alias or copy          & \hCode{Aliasing}        & \hCode{b = a} copies the arrow \\
			fill, accumulate       & \hCode{FillSquares}, \hCode{SumAverage} & \hCode{(double) sum / a.length} \\
			best so far            & \hCode{MinMax}          & start with index 0, loop from 1 \\
			search, count          & \hCode{LinearSearch}, \hCode{CountPassing} & stop at the first match; $-1$ \\
			transform, filter      & \hCode{CurveScores}, \hCode{FilterPassing} & count first, then copy \\
			swap, separators       & \hCode{Swap}, \hCode{Separators} & \hCode{temp}; a comma if \hCode{i > 0} \\
			partially filled       & \hCode{PartiallyFilled} & loops go to \hCode{size} \\
			methods                & \hCode{PassReference}   & elements yes, the caller's arrow no \\
			\bottomrule
		\end{tabular}
\end{table}%
% +++++++++++++++++++++++++++++++++++++++ A tbl


\end{frame}
% ~~~~~~~~~~~~~~~~~~~~~~~~~~~~~~~~~~~~~~ A frm




% ~~~~~~~~~~~~~~~~~~~~~~~~~~~~~~~~~~~~~~ V frm
\begin{frame}[t,fragile] \frametitle{Exercises: change the programs}
	\begin{enumerate}\small
		\item \hCode{AboveAverage}: also print how many days were \emph{below} the average.
		\item \hCode{MinMax}: make it report the \emph{last} of two equal maxima. Which character changes?
		\item \hCode{LinearSearch}: write \hCode{lastAtLeast}, which searches from the end.
		\item \hCode{FilterPassing}: make the limit a parameter of a method
		      \hCode{static int[] atLeast(int[] a, int limit)}.
		\item \hCode{Separators}: print \hCode{[72, 85, 90]} with the brackets, and \hCode{[]} for an empty array.
		\item \hCode{PartiallyFilled}: what happens if you add a ninth value? Add a check that prints
		      ``full'' instead.
		\item \hCode{PassReference}: write \hCode{reversed(int[] a)} that returns a new, reversed array.
	\end{enumerate}

	\vspace{0.5em}
	\small After a change, run \hCode{bash run\_examples.sh Name.java} in the deck folder
	(macOS: Terminal; Windows: Git Bash); it rewrites \hCode{codeoutput/Name.txt}.
	In Eclipse: copy the class into a project and \menu{Run As > Java Application}.
\end{frame}
% ~~~~~~~~~~~~~~~~~~~~~~~~~~~~~~~~~~~~~~ A frm




% ~~~~~~~~~~~~~~~~~~~~~~~~~~~~~~~~~~~~~~ V frm
\begin{frame}[t,fragile]
	\begin{center}
		\vspace{4cm}
		\hRemark{\Huge Questions?}
	\end{center}
\end{frame}
% ~~~~~~~~~~~~~~~~~~~~~~~~~~~~~~~~~~~~~~ A frm



\end{document}
